\section{학습, 평가와 모델 활성화}
\label{sec:learning}
\subsection{분할 단위와 데이터 역할}
학습 프로그램은 데이터의 역할과 분할 단위를 선언한다. 역할 집합을
$\{\mathsf{Fit},\mathsf{Tune},\mathsf{Eval},\mathsf{Adapt}\}$로 두고,
각 레코드에 피험자, 세션, 시행, 원본 구간 식별자를 부여한다.
분할 함수 $g$는 프로토콜에서 선택한 독립 단위를 반환한다.
\begin{equation}
 g:\mathrm{RecordID}\longrightarrow
 \mathrm{SubjectKey}\ \mid\ \mathrm{SessionKey}\ \mid\ \mathrm{TrialKey}.
\end{equation}
동일 원본 구간에서 생성한 창들은 공통 원천 식별자를 계보에 보관한다.
세션을 가로지르는 일반화 평가는 세션 key를,
피험자를 가로지르는 평가는 피험자 key를 분할 단위로 사용한다.
MOABB의 데이터셋ㆍ전처리ㆍ평가 구성을 이 데이터 역할의 연결 지점으로 사용한다\cite{moabb}.

\subsection{모델 객체와 fit 효과}
모델 버전 $v$를 다음 객체로 정의한다.
\begin{equation}
 M_v=(\vartheta_v,\mathcal C_{\mathrm{in}},\mathcal C_{\mathrm{out}},
      \rho_v^{\mathrm{fit}},a_v,\mathrm{digest}_v).
 \label{eq:model}
\end{equation}
$\vartheta_v$는 수치 매개변수, $\rho_v^{\mathrm{fit}}$는 매개변수와 선택 절차에
영향을 준 자료의 조상 집합, $a_v$는 해당 모델을 사용할 수 있게 된 시각이다.
공분산 기준, 정규화 통계, 특징 선택, 하이퍼파라미터 선택도 이 조상 집합에 포함한다.

학습 효과는 입력 자료의 역할을 소비하고 모델의 학습 계보를 생성한다.
고정 모델 평가의 적합 조건은 다음과 같다.
\begin{equation}
 \begin{aligned}
 \rho_v^{\mathrm{fit}}&\subseteq
 \Anc(D_{\mathrm{fit}}\cup D_{\mathrm{tune}}),\\
 g(\rho_v^{\mathrm{fit}})\cap g(D_{\mathrm{eval}})&=\varnothing.
 \end{aligned}
 \label{eq:split}
\end{equation}
컴파일러는 fit 효과가 생성한 모델 식별자를 추적하고, 평가 연산자는
식~\eqref{eq:split}과 입력 스키마를 검사한다.
이 조건은 필터ㆍ특징ㆍ추정기의 선택을 포함한 전체 학습 그래프에 적용한다.

중첩 평가에서는 외부 fold $j$마다 모델 $M^{(j)}$를 생성한다.
내부 fitㆍtune 자료는 외부 학습 집합의 부분집합으로 구성하고,
외부 평가 집합은 해당 fold의 결과 계산에 사용한다.
\begin{equation}
 \begin{aligned}
 D_{\mathrm{fit}}^{(j,k)}\cup D_{\mathrm{tune}}^{(j,k)}
 &\subseteq D_{\mathrm{train}}^{(j)},\\
 g(D_{\mathrm{train}}^{(j)})\cap g(D_{\mathrm{eval}}^{(j)})
 &=\varnothing.
 \end{aligned}
\end{equation}
분할 manifest는 원천 식별자, 그룹 함수, seed, 선택한 구간과 생성 절차의 버전을 보관한다.

\subsection{온라인 적응의 시간 의미}
적응 모델은 갱신 사건과 활성화 사건을 가진다. 버전 $v$의 활성화 시각을 $b_v$로 두면
\begin{equation}
 \max_{r\in\rho_v^{\mathrm{fit}}}a_r\le a_v\le b_v.
 \label{eq:modeltime}
\end{equation}
예측 작업은 시작 시점에 활성 모델 $M_{v(n)}$를 선택하고,
그 버전의 매개변수를 작업 완료까지 읽기 전용으로 보유한다.
\begin{equation}
 \widehat y_n=f_{\vartheta_{v(n)}}(x_n),\qquad
 b_{v(n)}\le a_n^{\mathrm{start}},\qquad
 \rho(\widehat y_n)=\rho(x_n)\cup\rho_{v(n)}^{\mathrm{fit}}.
\end{equation}
레이블을 사용하는 갱신은 레이블의 가용 시각을 추가 의존성으로 포함한다.
prequential 프로토콜은 예측, 결과 평가, 레이블 기반 갱신의 사건 순서를 기록한다.
비지도 적응은 입력 특징과 갱신 정책을 자신의 fit 계보에 기록한다.
정규화, 시계 보정과 calibration도 버전과 활성화 시각을 가지므로 같은 실행 규칙을 사용한다.

\subsection{추정값과 제어 사건}
추정기 출력은 class score, 확률 벡터, 연속 상태 추정량 등의 명목 형식을 갖는다.
확률 벡터는 $p_k\ge0$, $\sum_kp_k=1$을 만족하며,
확률 보정 연산은 방법ㆍ자료ㆍ버전을 추가한다.
의사결정 연산은 입력 형식과 임계값 정책을 사용하여 사건을 생성한다.
예를 들어
\begin{equation}
 d(p)=
 \begin{cases}
 \operatorname{Select}(\arg\max_kp_k),&\max_kp_k\ge\eta,\\
 \operatorname{Hold},&\max_kp_k<\eta
 \end{cases}
\end{equation}
로 정의한다. 최댓값의 동점은 선언한 클래스 순서의 첫 항목으로 결정한다.
임계값 $\eta$의 선택 자료는 해당 정책의 학습 계보에 포함한다.
명령 출력은 사건의 단위ㆍ범위ㆍ목표 시각과 장치 기능 요구를 입력받는다.
실행기는 출력한 명령과 장치 응답을 동일한 사건 식별자로 연결한다.

\subsection{평가 결과의 구조}
평가 연산자는 분할 manifest, 모델 버전, 예측ㆍ레이블 쌍, 품질 선택 정책과 지표 정의를 반환한다.
클래스별 recall을 평균하는 balanced accuracy는
\begin{equation}
 \mathrm{BA}=\frac1K\sum_{k=1}^{K}
 \frac{\mathrm{TP}_k}{\mathrm{TP}_k+\mathrm{FN}_k}
\end{equation}
로 계산하며, 각 분모가 양수인 평가 구성을 사용한다.
집계는 시행, 세션, 피험자 중 지정한 단위를 따라 수행한다.
구간 추정은 재표집 단위, 표본 수, 신뢰수준과 난수 seed를 결과 메타데이터로 포함한다.
지연 결과는 입력 준비, 출력 완료와 사건 응답의 관측 경계를 구분하여 기록한다.
